\chapter*{Quick Reference}
\addcontentsline{toc}{chapter}{Quick Reference}

This appendix provide a brief listing of the C language bindings of the
CDI library routines:

\section*{\tt \htmlref{gridDefNvertex}{gridDefNvertex}}

\begin{verbatim}
    void gridDefNvertex (int gridID, int nvertex);
\end{verbatim}

Define the number of vertex of a Gridbox.


\section*{\tt \htmlref{gridDefXbounds}{gridDefXbounds}}

\begin{verbatim}
    void gridDefXbounds (int gridID, const double xbounds[]);
\end{verbatim}

Define the bounds of a X-axis.


\section*{\tt \htmlref{gridDefYbounds}{gridDefYbounds}}

\begin{verbatim}
    void gridDefYbounds (int gridID, const double ybounds[]);
\end{verbatim}

Define the bounds of a Y-axis.


\section*{\tt \htmlref{gridInqNvertex}{gridInqNvertex}}

\begin{verbatim}
    int gridInqNvertex (int gridID);
\end{verbatim}

Get the number of vertex of a Gridbox.


\section*{\tt \htmlref{gridInqUUID}{gridInqUUID}}

\begin{verbatim}
    void gridInqUUID (int gridID, unsigned char uuid[CDI_UUID_SIZE]);
\end{verbatim}

Get the UUID of an unstructured grid.


\section*{\tt \htmlref{gridInqXbounds}{gridInqXbounds}}

\begin{verbatim}
    int gridInqXbounds (int gridID, double xbounds[]);
\end{verbatim}

Get the bounds of a X-axis.


\section*{\tt \htmlref{gridInqYbounds}{gridInqYbounds}}

\begin{verbatim}
    int gridInqYbounds (int gridID, double ybounds[]);
\end{verbatim}

Get the bounds of a Y-axis.


\section*{\tt \htmlref{subtypeCreate}{subtypeCreate}}

\begin{verbatim}
    int subtypeCreate (int subtype);
\end{verbatim}

Create a variable subtype.


\section*{\tt \htmlref{subtypeDefActiveIndex}{subtypeDefActiveIndex}}

\begin{verbatim}
    void subtypeDefActiveIndex (int subtypeID, int index);
\end{verbatim}

Set the currently active index of a subtype (e.g. current tile index)..


\section*{\tt \htmlref{subtypeInqActiveIndex}{subtypeInqActiveIndex}}

\begin{verbatim}
    int subtypeInqActiveIndex (int subtypeID);
\end{verbatim}

Get the currently active index of a subtype (e.g. current tile index)..


\section*{\tt \htmlref{subtypeInqSize}{subtypeInqSize}}

\begin{verbatim}
    int subtypeInqSize (int subtypeID);
\end{verbatim}

Get the size of a subtype (e.g. no. of tiles)..


\section*{\tt \htmlref{subtypeInqTile}{subtypeInqTile}}

\begin{verbatim}
    int subtypeInqTile (int subtypeID, int tileindex, int attribute);
\end{verbatim}

Specialized version of subtypeInqSubEntry looking for tile/attribute pair..


\section*{\tt \htmlref{tableRead}{tableRead}}

\begin{verbatim}
    int tableRead (const char *tablefile);
\end{verbatim}

read table of parameters from file in tabular format.


\section*{\tt \htmlref{tableWrite}{tableWrite}}

\begin{verbatim}
    void tableWrite (const char *filename, int tableID);
\end{verbatim}

write table of parameters to file in tabular format.


\section*{\tt \htmlref{tableWriteC}{tableWriteC}}

\begin{verbatim}
    void tableWriteC (const char *filename, int tableID);
\end{verbatim}

write table of parameters to file in C language format.


\section*{\tt \htmlref{taxisCreate}{taxisCreate}}

\begin{verbatim}
    int taxisCreate (int timetype);
\end{verbatim}

Create a Time axis.


\section*{\tt \htmlref{taxisDefCalendar}{taxisDefCalendar}}

\begin{verbatim}
    void taxisDefCalendar (int taxisID, int calendar);
\end{verbatim}

Define the calendar.


\section*{\tt \htmlref{taxisDefFdate}{taxisDefFdate}}

\begin{verbatim}
    void taxisDefFdate (int taxisID, int date);
\end{verbatim}

Define the forecast reference date.


\section*{\tt \htmlref{taxisDefFtime}{taxisDefFtime}}

\begin{verbatim}
    void taxisDefFtime (int taxisID, int time);
\end{verbatim}

Define the forecast reference time.


\section*{\tt \htmlref{taxisDefRdate}{taxisDefRdate}}

\begin{verbatim}
    void taxisDefRdate (int taxisID, int date);
\end{verbatim}

Define the reference date.


\section*{\tt \htmlref{taxisDefRtime}{taxisDefRtime}}

\begin{verbatim}
    void taxisDefRtime (int taxisID, int time);
\end{verbatim}

Define the reference time.


\section*{\tt \htmlref{taxisDefVdate}{taxisDefVdate}}

\begin{verbatim}
    void taxisDefVdate (int taxisID, int date);
\end{verbatim}

Define the verification date.


\section*{\tt \htmlref{taxisDefVtime}{taxisDefVtime}}

\begin{verbatim}
    void taxisDefVtime (int taxisID, int time);
\end{verbatim}

Define the verification time.


\section*{\tt \htmlref{taxisDestroy}{taxisDestroy}}

\begin{verbatim}
    void taxisDestroy (int taxisID);
\end{verbatim}

Destroy a Time axis.


\section*{\tt \htmlref{taxisInqCalendar}{taxisInqCalendar}}

\begin{verbatim}
    int taxisInqCalendar (int taxisID);
\end{verbatim}

Get the calendar.


\section*{\tt \htmlref{taxisInqFdate}{taxisInqFdate}}

\begin{verbatim}
    int taxisInqFdate (int taxisID);
\end{verbatim}

Get the forecast reference date.


\section*{\tt \htmlref{taxisInqFtime}{taxisInqFtime}}

\begin{verbatim}
    int taxisInqFtime (int taxisID);
\end{verbatim}

Get the forecast reference time.


\section*{\tt \htmlref{taxisInqRdate}{taxisInqRdate}}

\begin{verbatim}
    int taxisInqRdate (int taxisID);
\end{verbatim}

Get the reference date.


\section*{\tt \htmlref{taxisInqRtime}{taxisInqRtime}}

\begin{verbatim}
    int taxisInqRtime (int taxisID);
\end{verbatim}

Get the reference time.


\section*{\tt \htmlref{taxisInqVdate}{taxisInqVdate}}

\begin{verbatim}
    int taxisInqVdate (int taxisID);
\end{verbatim}

Get the verification date.


\section*{\tt \htmlref{taxisInqVtime}{taxisInqVtime}}

\begin{verbatim}
    int taxisInqVtime (int taxisID);
\end{verbatim}

Get the verification time.


\section*{\tt \htmlref{vlistInqVarSubtype}{vlistInqVarSubtype}}

\begin{verbatim}
    int vlistInqVarSubtype (int vlistID, int varID);
\end{verbatim}

Return subtype ID for a given variable..


\section*{\tt \htmlref{zaxisCreate}{zaxisCreate}}

\begin{verbatim}
    int zaxisCreate (int zaxistype, int size);
\end{verbatim}

Create a vertical Z-axis.


\section*{\tt \htmlref{zaxisDefLevel}{zaxisDefLevel}}

\begin{verbatim}
    void zaxisDefLevel (int zaxisID, int levelID, double levels);
\end{verbatim}

Define one level of a Z-axis.


\section*{\tt \htmlref{zaxisDefLevels}{zaxisDefLevels}}

\begin{verbatim}
    void zaxisDefLevels (int zaxisID, const double levels[]);
\end{verbatim}

Define the levels of a Z-axis.


\section*{\tt \htmlref{zaxisDefLongname}{zaxisDefLongname}}

\begin{verbatim}
    void zaxisDefLongname (int zaxisID, const char *longname_optional);
\end{verbatim}

Define the longname of a Z-axis.


\section*{\tt \htmlref{zaxisDefName}{zaxisDefName}}

\begin{verbatim}
    void zaxisDefName (int zaxisID, const char *name_optional);
\end{verbatim}

Define the name of a Z-axis.


\section*{\tt \htmlref{zaxisDefNlevRef}{zaxisDefNlevRef}}

\begin{verbatim}
    void zaxisDefNlevRef (int gridID, int nhlev);
\end{verbatim}

Define the number of half levels of a generalized Z-axis.


\section*{\tt \htmlref{zaxisDefNumber}{zaxisDefNumber}}

\begin{verbatim}
    void zaxisDefNumber (int gridID, int number);
\end{verbatim}

Define the reference number for a generalized Z-axis.


\section*{\tt \htmlref{zaxisDefPsName}{zaxisDefPsName}}

\begin{verbatim}
    void zaxisDefPsName (int zaxisID, const char *psname_optional);
\end{verbatim}

Define the name of the surface pressure variable of a hybrid sigma pressure Z-axis.


\section*{\tt \htmlref{zaxisDefUUID}{zaxisDefUUID}}

\begin{verbatim}
    void zaxisDefUUID (int zaxisID, const unsigned char uuid[CDI_UUID_SIZE]);
\end{verbatim}

Define the UUID of a generalized Z-axis.


\section*{\tt \htmlref{zaxisDefUnits}{zaxisDefUnits}}

\begin{verbatim}
    void zaxisDefUnits (int zaxisID, const char *units_optional);
\end{verbatim}

Define the units of a Z-axis.


\section*{\tt \htmlref{zaxisDestroy}{zaxisDestroy}}

\begin{verbatim}
    void zaxisDestroy (int zaxisID);
\end{verbatim}

Destroy a vertical Z-axis.


\section*{\tt \htmlref{zaxisDuplicate}{zaxisDuplicate}}

\begin{verbatim}
    int zaxisDuplicate (int zaxisID);
\end{verbatim}

Duplicate a Z-axis.


\section*{\tt \htmlref{zaxisInqLevel}{zaxisInqLevel}}

\begin{verbatim}
    double zaxisInqLevel (int zaxisID, int levelID);
\end{verbatim}

Get one level of a Z-axis.


\section*{\tt \htmlref{zaxisInqLevels}{zaxisInqLevels}}

\begin{verbatim}
    void zaxisInqLevels (int zaxisID, double levels[]);
\end{verbatim}

Get all levels of a Z-axis.


\section*{\tt \htmlref{zaxisInqLongname}{zaxisInqLongname}}

\begin{verbatim}
    void zaxisInqLongname (int zaxisID, char *longname);
\end{verbatim}

Get the longname of a Z-axis.


\section*{\tt \htmlref{zaxisInqName}{zaxisInqName}}

\begin{verbatim}
    void zaxisInqName (int zaxisID, char *name);
\end{verbatim}

Get the name of a Z-axis.


\section*{\tt \htmlref{zaxisInqNlevRef}{zaxisInqNlevRef}}

\begin{verbatim}
    int zaxisInqNlevRef (int gridID);
\end{verbatim}

Get the number of half levels of a generalized Z-axis.


\section*{\tt \htmlref{zaxisInqNumber}{zaxisInqNumber}}

\begin{verbatim}
    int zaxisInqNumber (int gridID);
\end{verbatim}

Get the reference number to a generalized Z-axis.


\section*{\tt \htmlref{zaxisInqPsName}{zaxisInqPsName}}

\begin{verbatim}
    void zaxisInqPsName (int zaxisID, char *psname);
\end{verbatim}

Get the name of the surface pressure variable of a hybrid sigma pressure Z-axis.


\section*{\tt \htmlref{zaxisInqSize}{zaxisInqSize}}

\begin{verbatim}
    int zaxisInqSize (int zaxisID);
\end{verbatim}

Get the size of a Z-axis.


\section*{\tt \htmlref{zaxisInqStdname}{zaxisInqStdname}}

\begin{verbatim}
    void zaxisInqStdname (int zaxisID, char *stdname);
\end{verbatim}

Get the standard name of a Z-axis.


\section*{\tt \htmlref{zaxisInqType}{zaxisInqType}}

\begin{verbatim}
    int zaxisInqType (int zaxisID);
\end{verbatim}

Get the type of a Z-axis.


\section*{\tt \htmlref{zaxisInqUUID}{zaxisInqUUID}}

\begin{verbatim}
    void zaxisInqUUID (int zaxisID, unsigned char uuid[CDI_UUID_SIZE]);
\end{verbatim}

Get the UUID of a generalized Z-axis.


\section*{\tt \htmlref{zaxisInqUnits}{zaxisInqUnits}}

\begin{verbatim}
    void zaxisInqUnits (int zaxisID, char *units);
\end{verbatim}

Get the units of a Z-axis.


